\documentclass[11pt,a4paper]{article}

\usepackage[margin=1in]{geometry}
\usepackage{hyperref,amsmath,amssymb,graphicx,enumitem}
\hypersetup{colorlinks=true,linkcolor=blue,urlcolor=blue}

\title{\textbf{Compiled Language Models}\\[4pt]
       \large Zero-Gradient Text Generation and the Path Forward}
\author{Douglas Rawson\\ \texttt{rawson.douglas@gmail.com}}
\date{May 2026}

\begin{document}
\maketitle

\begin{abstract}
We present a working compiled language model---a transformer that generates text using only
deterministically constructed weight matrices. No gradient descent. No training. Every parameter
is specified by closed-form computation from corpus statistics.

On the TinyStories corpus, our compiled LM achieves 35.71\% held-out top-1 next-token accuracy
using 716,000 compiled rules across a 4-gram, trigram, bigram, and macro backoff hierarchy with
an attention-based position shift mechanism. On WikiText-103, a 500,000-rule trigram-only model
achieves 2.44\%---98$\times$ better than random but far below the TinyStories result.

We characterize the gap: it is architectural, not fundamental. The TinyStories LM uses
multi-order backoff and attention; the Wikipedia LM uses a single trigram layer. We outline the
specific engineering steps needed to close this gap and invite collaboration from researchers
in mechanistic interpretability, efficient architectures, and language modeling.
\end{abstract}

\section{Introduction}

The prevailing approach to building language models is to train them. Billions of parameters.
Trillions of tokens. Thousands of GPU-hours. The weights that result are opaque---we know they
work, but we don't know \emph{which} parameters encode \emph{which} linguistic knowledge.

We take a different approach: \textbf{compile} the weights directly. Every parameter in our model
is computed by a deterministic formula from corpus statistics. Every rule in the FFN can be
inspected, audited, and traced back to its source n-gram. The model never sees a gradient.

This paper describes what works, what doesn't, and what needs to happen next.

\section{How It Works}

\subsection{Architecture}

The compiled LM uses a transformer-style architecture with hand-constructed weights:

\begin{itemize}[leftmargin=*]
\item \textbf{Embeddings}: PPMI+SVD from corpus co-occurrence statistics (Levy \& Goldberg, 2014). No random initialization, no training. 8,000-word vocabulary for TinyStories; 4,000-token BPE vocabulary for Wikipedia.
\item \textbf{Position encoding}: An orthonormal basis $\{p_0, \ldots, p_{T-1}\}$ of dimension $T$ where $p_i \cdot p_j = \delta_{ij}$. A closed-form shift operator $P_1 = \sum_{i \geq 1} p_{i-1} \otimes p_i$ provides exact ``attend to previous position'' without learned positional embeddings.
\item \textbf{Attention}: Two hand-written attention heads. Head 1 (previous-token head): $Q = P_1$, $K = I$, $V = I$, writes into a dedicated PREV residual slot. Head 2 (induction head): $Q$ reads the current token embedding, $K$ reads PREV, $V$ reads the main residual---the classic Olsson induction circuit, realized exactly through matrix construction.
\item \textbf{FFN rules}: The core language modeling capability. Each compiled rule is a rank-1 FFN entry: a trigger embedding direction (the n-gram context), a threshold, and a response embedding direction (the predicted next token). The TinyStories LM uses 500K 4-gram rules + 200K trigram + 16K bigram + 10K macros, organized in a confidence-gated backoff hierarchy.
\item \textbf{Unembedding}: Tied to the embedding matrix. No separate output layer training.
\end{itemize}

\subsection{The Forward Pass}

For a history of token IDs $\{t_1, \ldots, t_n\}$, the compiled LM:

\begin{enumerate}[leftmargin=*]
\item Embeds each token into a residual stream $h \in \mathbb{R}^d$.
\item Applies the previous-token attention head to populate a PREV slot.
\item Applies the induction head, which uses the current token (Q) and previous token (K) to attend to past occurrences of the same pattern.
\item Runs the FFN rule cascade: attempts 4-gram lookup. If confidence $>$ threshold, uses it. Otherwise falls back to trigram, then bigram, then macros. If no rule fires, returns uniform logits.
\item Projects the residual through the tied embedding matrix to produce vocabulary logits.
\end{enumerate}

All of these operations are linear algebra with hand-specified matrices. There is no softmax temperature tuning, no learned scaling factors---every constant was derived analytically.

\section{Results}

\subsection{TinyStories (8,000-word vocabulary)}

Evaluated on 24,581 held-out positions from 200 held-out stories (stories 500,000--500,200, never used during compilation).

\begin{table}[h]
\centering
\begin{tabular}{lrr}
\hline
\textbf{Metric} & \textbf{Random Embed} & \textbf{PPMI Embed} \\
\hline
Top-1 accuracy (in-vocab) & \textbf{35.71\%} & 32.69\% \\
Top-5 accuracy (in-vocab) & \textbf{54.97\%} & 47.63\% \\
Perplexity (in-vocab) & 2.1M & 5.4M \\
Fire rate & 100\% & 100\% \\
\hline
\end{tabular}
\caption{Held-out evaluation on TinyStories. PPMI embeddings improve cosine-tolerant accuracy (synonyms count as correct) but hurt exact-match due to semantic clustering of predictions.}
\end{table}

The random-embedding model achieves 35.71\% top-1 accuracy---predicting the next word correctly
more than one third of the time on never-seen text, using only compiled rules. The PPMI model's
lower exact-match accuracy is a known metric artifact: PPMI predictions cluster around semantic
neighbors of the target (e.g., predicting ``mum'' when the target is ``mom''), which exact-match
evaluation penalizes. Under a cosine-tolerant metric ($\cos > 0.5$), the PPMI gap shrinks to
$-0.54$ pp (from $-3.03$ pp exact).

\subsection{WikiText-103 (4,000-token BPE vocabulary)}

Evaluated on 19,113 held-out positions from 200 held-out paragraphs (paragraphs 700,000--700,199).

\begin{table}[h]
\centering
\begin{tabular}{lrr}
\hline
\textbf{Metric} & \textbf{All Positions} & \textbf{In-Vocab (84.8\%)} \\
\hline
Top-1 accuracy & 2.07\% & 2.44\% \\
Top-5 accuracy & 4.11\% & 4.85\% \\
Perplexity (in-vocab) & --- & 3,645 \\
Uniform baseline & --- & 4,000 \\
Fire rate & 100\% & 100\% \\
\hline
\end{tabular}
\caption{Held-out evaluation on WikiText-103. The model uses a single-order trigram architecture with 500,000 rules. Results improve marginally with larger context (128 vs 64 tokens) and more active rules (32 vs 8), suggesting headroom.}
\end{table}

The Wikipedia LM is 98$\times$ better than random (2.44\% vs 0.025\%) but dramatically below
the TinyStories result. Three factors explain the gap:

\begin{enumerate}[leftmargin=*]
\item \textbf{Architecture}: Single-order trigram lookup (rule: ``given two previous tokens, predict next'') vs. multi-order backoff (4-gram $\rightarrow$ trigram $\rightarrow$ bigram $\rightarrow$ macro). The backoff hierarchy quadruples the effective context and gracefully handles unseen n-grams.
\item \textbf{Attention}: The TinyStories LM uses a compiled induction head that copies patterns from earlier in the text. The Wikipedia LM uses only mean-pooled embeddings of the history window, losing positional information.
\item \textbf{Task difficulty}: TinyStories is a constrained domain of children's stories with predictable narrative patterns. Wikipedia spans every domain of human knowledge with highly variable text structure.
\end{enumerate}

\section{What This Means}

The compiled LM is not a curiosity. 35.71\% held-out accuracy on real text, from a model whose
every parameter was written by a Python script---no training, no gradient descent, no
hyperparameter search---demonstrates that \emph{language modeling can be compiled}.

The architecture has clear headroom. The gap between 35\% (TinyStories) and 2.4\% (Wikipedia) is
not a ceiling---it is the difference between a production architecture and a simplified one.
Closing it requires \emph{engineering}, not breakthrough research:

\begin{enumerate}[leftmargin=*]
\item \textbf{Multi-order backoff for Wikipedia}: Port the 4-gram+trigram+bigram+macro hierarchy from the TinyStories compiler to the Wikipedia BPE tokenizer. Expected impact: 5--10$\times$ improvement.
\item \textbf{Proper threshold calibration}: The TinyStories model uses per-rule thresholds based on off-target cosine similarity. The Wikipedia model uses a uniform zero threshold---every rule fires for every query. Calibrating thresholds would reduce noise and improve precision.
\item \textbf{Compiled induction for BPE text}: The induction head enables pattern copying across positions. Wiring it to the BPE tokenizer would capture long-range dependencies.
\item \textbf{PPMI quality for larger corpora}: The Wikipedia PPMI embeddings use a 10\% subsample of trigrams. A full-count PPMI with proper SVD would improve the embedding geometry.
\end{enumerate}

\section{Related Work}

The compiled approach draws from several traditions:

\begin{itemize}[leftmargin=*]
\item \textbf{Transformer circuits} (Elhage et al., 2021; Olsson et al., 2022): The compiled induction head is an exact realization of the attention pattern identified in trained transformers. We show it can be constructed analytically rather than discovered through interpretability.
\item \textbf{n-gram language models}: The FFN rule system is fundamentally an n-gram model embedded in a transformer architecture. The backoff hierarchy is the standard Katz backoff, implemented as a confidence-gated FFN cascade rather than a statistical model.
\item \textbf{Word embeddings} (Levy \& Goldberg, 2014): The PPMI+SVD embedding computation is the classical word2vec algorithm executed in closed form. No SGD required.
\item \textbf{Model editing} (Meng et al., 2022): Techniques like ROME edit specific facts in trained models. The compiled LM \emph{starts} from the edited state---every fact is placed deliberately.
\end{itemize}

\section{Call for Collaboration}

We have demonstrated that language models can be compiled. The path to competitive performance
is clear but labor-intensive. We invite collaboration on:

\begin{itemize}[leftmargin=*]
\item \textbf{Multi-order compiler for Wikipedia}: Porting the TinyStories v2 compiler pipeline to BPE-tokenized corpora with the full rule hierarchy. Python and PyTorch.
\item \textbf{Efficient n-gram extraction}: The current pipeline uses 500K--716K rules. Scaling to tens of millions of rules requires efficient sparse matrix operations and potentially approximate nearest-neighbor search.
\item \textbf{Threshold calibration}: Systematic study of per-rule threshold settings for compiled FFN architectures, including noise budget analysis for different corpus sizes.
\item \textbf{PPMI at scale}: Full-count PPMI on billion-token corpora. The current Wikipedia PPMI uses a 10\% subsample due to memory constraints. Efficient sparse SVD implementations (e.g., randomized SVD) would help.
\item \textbf{Compiled architectures for code}: Beyond natural language, compiled modules for code completion, API recall, and error correction.
\item \textbf{Safety and auditability}: The compiled LM's complete parameter auditability makes it uniquely suited for safety research. Every prediction can be traced to specific rules. We welcome red-teaming and verification research.
\end{itemize}

\section{Conclusion}

Compiled language models are real. They work. They are far from competitive with trained models,
but they are built on a fundamentally different principle: every parameter has a known purpose,
every rule has a source in the training corpus, and no optimization was required to arrive at
any of them.

The 35.71\% result on TinyStories is a proof of concept. The 2.44\% result on Wikipedia is a
baseline. The engineering work to close the gap is defined and tractable. We invite the
community to help build it.

\vspace{1em}
\noindent\textit{Repository: \url{https://github.com/DRawson5570/compiled-injection}}\\
\textit{Contact: \texttt{rawson.douglas@gmail.com}}\\
\textit{All experiments documented in \texttt{docs/EXPERIMENT\_LOG.md}.}\\
\textit{Provenance: OpenTimestamp + Zenodo DOI.}

\end{document}
